% Part: computability
% Chapter: computability-theory
% Section: normal-form

\documentclass[../../../include/open-logic-section]{subfiles}

\begin{document}

\olfileid[hi]{cmp}{thy}{nfm}
\olsection{प्रसामान्य रूप प्रमेय}

मान लें कि हम संगणनीय फलनों की परिभाषाओं का वर्णन कर सकते हैं और यह जाँच
सकते हैं कि किसी आगत पर उस फलन की संगणना के पूर्ण अभिलेख का कोई प्रस्तावित
वर्णन सही है या नहीं। तब यह उचित ही है कि संगणना के मॉडल से स्वतंत्र होकर
हम किसी भी संगणनीय फलन~$f$ का किसी भी आगत~$x$ पर मान इस प्रकार निर्धारित
कर सकते हैं:
\begin{enumerate}
  \item संगणना-अभिलेखों के सभी संभव वर्णनों में खोज करें।
  \item जाँचें कि कोई दिया हुआ अभिलेख~$f(x)$ की संगणना का अभिलेख है या नहीं।
  \item सही अभिलेख मिल जाने पर उससे~$f(x)$ का मान निकालें।
\end{enumerate}
वास्तव में ऐसा ही होता है—यही क्लीन के प्रसामान्य रूप प्रमेय का विषय है।

\begin{thm}[क्लीन का प्रसामान्य रूप प्रमेय]
\ollabel{thm:normal-form}
एक आदिम पुनरावर्ती संबंध~$T(e, x, s)$ और एक आदिम पुनरावर्ती फलन~$U(s)$
हैं जिनमें निम्न गुण है: यदि $f$ कोई आंशिक संगणनीय फलन है, तो किसी~$e$ के लिए,
\[
f(x) \simeq U(\umin{s}{T(e, x, s)})
\]
प्रत्येक~$x$ के लिए।
\end{thm}

\begin{proof}[प्रमाण की रूपरेखा]
संगणना के किसी भी मॉडल के लिए संगणनीय फलन~$f$ के वर्णन को कठोरता से
परिभाषित किया जा सकता है और ऐसे वर्णन को किसी प्राकृतिक संख्या~$e$ से
कूटित किया जा सकता है। इसी प्रकार ``संगणना अनुक्रम'' की धारणा को कठोरता
से परिभाषित किया जा सकता है, जो सूचकांक~$e$ वाले फलन की आगत~$x$ पर
संगणना की प्रक्रिया को अभिलिखित करती है। ऐसे संगणना अनुक्रम को भी किसी
संख्या~$s$ से कूटित किया जा सकता है। यह इस प्रकार किया जा सकता है कि
\begin{enumerate}
  \item संबंध $T(e, x, s)$, जो तभी सत्य होता है जब संख्या~$s$, सूचकांक~$e$
  वाले फलन के आगत~$x$ पर संगणना अनुक्रम को कूटित करती है, और
  \item फलन $U(s)$, जो~$s$ से कूटित संगणना अनुक्रम को उस संगणना के
  अंतिम परिणाम पर प्रतिचित्रित करता है,
\end{enumerate}
दोनों संगणनीय हैं। वास्तव में संबंध~$T$ और फलन~$U$ आदिम पुनरावर्ती हैं।
\end{proof}

\begin{explain}
प्रसामान्य रूप प्रमेय का कठोर प्रमाण देने के लिए हमें संगणना का कोई मॉडल
नियत करना होगा, संगणनीय फलनों के वर्णनों और संगणना अनुक्रमों का कूटांकन
विस्तार से करना होगा, और सत्यापित करना होगा कि संबंध~$T$ तथा फलन~$U$
आदिम पुनरावर्ती हैं। अधिकांश अनुप्रयोगों के लिए इतना पर्याप्त है कि $T$
और~$U$ संगणनीय हों तथा $U$~पूर्ण हो।

प्रसामान्य रूप प्रमेय के प्रमाण को इस सूत्र-वाक्य में याद रखना संभवतः
सबसे अच्छा है: $\umin{s}{T(e, x, s)}$, सूचकांक~$e$ वाले फलन की आगत~$x$
पर संगणना के अनुक्रम की खोज करता है, और ऐसा अनुक्रम मिल जाने पर $U$
उस संगणना अनुक्रम का निर्गत लौटाता है।
\end{explain}

यदि संगणना का मॉडल आंशिक पुनरावर्ती फलनों का है (जिसका क्लीन ने मूलतः
उपयोग किया था), तो यह दिखाता है कि किसी भी फलन की परिभाषा में अपरिबद्ध
खोज का केवल एक प्रयोग, अर्थात् केवल एक $\umin{y}{f(\vec x, y)}$ संकारक
आवश्यक है। इस अर्थ में यह दिखाता है कि प्रत्येक आंशिक पुनरावर्ती फलन का
एक प्रसामान्य रूप होता है।

$T$ और $U$ का उपयोग परिगणन $\cfind{0}$, $\cfind{1}$, $\cfind{2}$,
\dots को परिभाषित करने में किया जा सकता है। अब से हम मानेंगे कि हमने
$T$ और~$U$ का कोई उपयुक्त चयन नियत कर लिया है, और समीकरण
\[
\cfind{e}(x) \simeq U(\umin{s}{T(e,x,s)})
\]
को $\cfind{e}$ की \emph{परिभाषा} मानेंगे।

एक और उपयोगी तथ्य यह है:

\begin{thm}
प्रत्येक आंशिक संगणनीय फलन के अपरिमित रूप से अनेक सूचकांक होते हैं।
\end{thm}

यह भी सहज रूप से स्पष्ट है। किसी संगणनीय फलन (के वर्णन) से आरंभ करके
हम एक अलग वर्णन बना सकते हैं जो उसी फलन (आगत–निर्गत युग्म) को संगणित
करता है, किंतु पहले ऐसा कोई कार्य करता है जिसका संगणना पर कोई प्रभाव
नहीं पड़ता (मसलन, जाँचता है कि $0 = 0$, अथवा $5$ तक गिनता है, आदि)।
परिवर्तित वर्णन का सूचकांक सदैव मूल सूचकांक से भिन्न होगा। दोनों उसी फलन
के सूचकांक हैं; बस उसकी संगणना कुछ भिन्न ढंग से की गई है।

\end{document}
